\chapter{Implementacja}

\section{Środowisko implementacyjne}
Implementację przygotowano jako zestaw artefaktów \emph{Infrastructure as Code}, z podziałem na:
\begin{itemize}
    \item warstwę Terraform (\texttt{providers.tf}, \texttt{variables.tf}, \texttt{main.tf}, \texttt{outputs.tf}),
    \item konfigurację Istio (\texttt{istiod-values.yaml}, \texttt{namespace-5gc.yaml}),
    \item manifesty workloadów i danych (\texttt{manifests/*}).
\end{itemize}

Podejście to umożliwia odtwarzalne wdrożenia oraz formalną kontrolę zmian w repozytorium kodu.

\section{Implementacja warstwy platformowej}
W Terraform zdefiniowano parametry pojemnościowe klastra dla obciążenia 5G SBA, w tym:
\begin{itemize}
    \item pule workerów dla ścieżki produkcyjnej i shadow,
    \item namespace \texttt{5gc} z etykietą \texttt{istio-injection=enabled},
    \item limitowanie zasobów poprzez \texttt{ResourceQuota}.
\end{itemize}

Przykładowy fragment deklaracji namespace:
\begin{lstlisting}[language=JavaScript,style=JavaScriptStyle,caption=Definicja namespace z auto-injection,label=lst:ns-5gc]
resource "kubernetes_namespace" "fivegc" {
  metadata {
    name = var.fivegc_namespace
    labels = {
      "istio-injection" = "enabled"
      "platform"        = "5g-core"
      "shadow-mode"     = "enabled"
    }
  }
}
\end{lstlisting}

\section{Implementacja izolacji danych}
W warstwie danych przygotowano dwa niezależne StatefulSety:
\begin{itemize}
    \item \texttt{udr-prod} -- baza produkcyjna,
    \item \texttt{udr-shadow} -- baza testowa dla wersji kandydackiej.
\end{itemize}

Do każdej bazy przypisano oddzielny sekret i osobną usługę Kubernetes. Rozdzielenie poświadczeń uniemożliwia wykorzystanie danych produkcyjnych przez SMF v2.

\section{Implementacja SMF v1 i SMF v2}
Wersja stabilna \texttt{smf-v1} i wersja shadow \texttt{smf-v2} zostały uruchomione jako niezależne deploymenty. Obie wersje są selekcjonowane przez etykietę \texttt{version}, co pozwala na przypisanie ich do subsetów Istio.

Izolację połączeń bazodanowych zrealizowano przez oddzielne \texttt{ConfigMap}:
\begin{itemize}
    \item \texttt{smf-v1-config} wskazuje \texttt{udr-prod},
    \item \texttt{smf-v2-config} wskazuje \texttt{udr-shadow}.
\end{itemize}

\section{Implementacja traffic mirroring}
Mechanizm mirrorowania zdefiniowano przez \texttt{DestinationRule} i \texttt{VirtualService}. Konfiguracja zapewnia:
\begin{itemize}
    \item obsługę całego ruchu przez \texttt{v1},
    \item przesył kopii 10\% żądań do \texttt{v2}.
\end{itemize}

Przykładowy fragment reguły:
\begin{lstlisting}[language=JavaScript,style=JavaScriptStyle,caption=Fragment VirtualService z mirrorowaniem,label=lst:vs-mirror]
http:
  - name: smf-shadow-mirroring
    route:
      - destination:
          host: smf.5gc.svc.cluster.local
          subset: v1
        weight: 100
    mirror:
      host: smf.5gc.svc.cluster.local
      subset: v2
    mirrorPercentage:
      value: 10.0
\end{lstlisting}

\section{Implementacja walidacji jakości}
Warstwę walidacyjną zbudowano na Iter8 i Prometheus. Eksperyment Iter8:
\begin{itemize}
    \item pobiera metryki \texttt{latency\_avg\_ms} i \texttt{error\_rate} zapytaniami PromQL,
    \item porównuje wynik z progami SLO,
    \item zwraca werdykt \texttt{pass/fail}.
\end{itemize}

Dodatkowo zdefiniowano opcjonalny deployment Diffy do porównania odpowiedzi API wersji stabilnej i kandydackiej z pominięciem pól zaszumionych (np. znaczniki czasu).

